% القسم: المناهج
% الفصل: الاستقراء
% المبحث: الاستقراء البنيوي

\documentclass[../../../include/open-logic-section]{subfiles}

\begin{document}

\olfileid{mth}{ind}{sti}

\olsection{الاستقراء البنيوي}

جرى استعمال الاستقراء إلى الآن في أحكام تعم الأعداد الطبيعية. وله
نظير مباشر يثبت الأحكام لجميع !!{element} مجموعة معرّفة استقرائيًّا.
ويسمى هذا \emph{الاستقراء البنيوي}، لأن الحجة فيه تتبع بنية الأشياء
التي أنشأها التعريف الاستقرائي.

وقوام التعريف الاستقرائي أمران: (أ)~تعيين !!{element} «الأولية» في
المجموعة، و(ب)~تعيين العمليات التي تنشئ !!{element} جديدة منها.
فالأشياء الأولية في مثال الحدود السليمة هي الحروف، والعملية الوحيدة
هي:
\begin{align*}
  o(s_1, s_2) &= [s_1 \circ s_2]
\end{align*}
ويجوز أن تُرى الأعداد الطبيعية~$\Nat$ نفسها ثمرة تعريف استقرائي:
فالابتداء بـ~$0$، والعملية دالة الخلف~$x + 1$.

وإذا أُريد إثبات حكم في جميع أفراد مجموعة معرّفة استقرائيًّا، أي
إثبات أن كل !!{element} فيها ذو خاصية~$P$، وجب أمران:
\begin{enumerate}
\item إثبات الخاصية~$P$ للأشياء الأولية.
\item وإثبات، لكل عملية~$o$، أن اتصاف معاملاتها بـ~$P$ يوجب اتصاف
  نتيجتها بها أيضًا.
\end{enumerate}
فالحكم على جميع الحدود السليمة، مثلًا، يثبت أولًا لكل حرف، ثم يثبت
لـ$[s_1 \circ s_2]$ متى ثبت لكل من $s_1$ و$s_2$ منفردًا.

\begin{prop}
  في كل حد سليم~$t$، يساوي عدد الرموز $[$ عدد الرموز $]$.
\end{prop}

\begin{proof}
نستعمل الاستقراء البنيوي؛ فالحدود السليمة معرّفة استقرائيًّا، والحروف
أفرادها الأولية، و$o$ هي العملية المنشئة للحد الجديد من حدين سابقين.
\begin{enumerate}
\item القضية صادقة لكل حرف؛ فعدد الرموز $[$ في الحرف المنفرد~$0$،
  وعدد الرموز $]$ فيه~$0$ أيضًا.
\item افرض أن عدد الرموز $[$ في $s_1$ يساوي عدد الرموز $]$ في
  $s_1$، وأن النظير صادق لـ$s_2$. فعدد الرموز $[$ في
  $o(s_1, s_2)$، أي في $[s_1 \circ s_2]$، هو مجموع عدديها في
  $s_1$ و$s_2$ مع زيادة واحد. وعدد الرموز $]$ في
  $o(s_1, s_2)$ هو كذلك مجموع عدديها في $s_1$ و$s_2$ مع زيادة
  واحد. ومن ثم يتساوى عدد الرموز $[$ في $o(s_1, s_2)$ وعدد الرموز
  $]$ فيه.
\end{enumerate}
\end{proof}

\begin{prob}
  أثبت بالاستقراء البنيوي أنه لا يبدأ أي حد سليم بالرمز~$]$.
\end{prob}

ولنأت ببرهان بنيوي آخر. القطعة الابتدائية الصحيحة من سلسلة الرموز~$t$
هي سلسلة~$s$ تطابق $t$ رمزًا رمزًا من جهة اليسار، على أن يكون $t$
أطول منها. فـ$[a \circ {}$، مثلًا، قطعة ابتدائية صحيحة من
$[a \circ b]$؛ وليس $[b \circ {}$ كذلك لاختلاف الرمز الثاني، ولا
$[a \circ b]$ نفسها لتساوي الطولين.

\begin{prop}\ollabel{prop:initial}
  في كل قطعة ابتدائية صحيحة غير خالية من حد سليم~$t$ عدد من الرموز $[$ أكبر من
  عدد الرموز $]$.
\end{prop}

\begin{proof}
  نستعمل الاستقراء على~$t$:
  \begin{enumerate}
  \item $t$ حرف بمفرده: عندئذ ليست لـ$t$ أي قطعة ابتدائية صحيحة.
  \item $t = [s_1 \circ s_2]$ لبعض حدين سليمين $s_1$ و~$s_2$. إذا
    كان $r$ قطعة ابتدائية صحيحة من $t$، فثمة عدة احتمالات:
    \begin{enumerate}
    \item $r$ هي $[$ فقط: عندئذ في $r$ رمز $[$ واحد أكثر من الرموز
      $]$.
    \item $r$ هي $[r_1$، حيث $r_1$ قطعة ابتدائية صحيحة من~$s_1$:
      بما أن $s_1$ حد سليم، فإن $r_1$، بحسب فرض الاستقراء، تحتوي على
      رموز $[$ أكثر من الرموز $]$؛ وكذلك $[r_1$.
    \item $r$ هي $[s_1$ أو $[s_1 \circ {}$: بحسب النتيجة السابقة،
      يتساوى عدد الرموز $[$ و$]$ في~$s_1$؛ ولذلك يزيد عدد الرموز $[$
      في $[s_1$ أو $[s_1 \circ {}$ بواحد على عدد الرموز~$]$.
    \item $r$ هي $[s_1 \circ r_2$، حيث $r_2$ قطعة ابتدائية صحيحة
      من~$s_2$: بحسب فرض الاستقراء، تحتوي $r_2$ على رموز $[$ أكثر من
      الرموز $]$. وبحسب النتيجة السابقة، يتساوى عدد الرموز $[$ و$]$
      في~$s_1$. لذلك يكون عدد الرموز $[$ في $[s_1 \circ r_2$ أكبر من
      عدد الرموز~$]$.
    \item $r$ هي $[s_1 \circ s_2$: بحسب النتيجة السابقة، يتساوى عدد
      الرموز $[$ و$]$ في $s_1$، ويصدق الشيء نفسه على~$s_2$. ولذلك
      يوجد في $[s_1 \circ s_2$ رمز $[$ واحد أكثر من الرموز~$]$.
    \end{enumerate}
  \end{enumerate}
\end{proof}

\end{document}
